\documentclass[11pt]{scrartcl}

\usepackage[top=1.5cm]{geometry}
\usepackage{url}
\usepackage{float}
\usepackage{listings}
\usepackage{xcolor}
\usepackage{graphicx}
\usepackage{booktabs}

\setlength{\parindent}{0em}
\setlength{\parskip}{0.5em}

\newcommand{\youranswerhere}{[Your answer goes here \ldots]}
\renewcommand{\thesubsection}{\arabic{subsection}}

\lstdefinestyle{dmrsql}{
  language=SQL,
  basicstyle=\small\ttfamily,
  keywordstyle=\color{magenta!75!black},
  stringstyle=\color{green!50!black},
  showspaces=false,
  showstringspaces=false,
  commentstyle=\color{gray}}

\lstdefinestyle{dmrJava}{
  language=JAVA,
  basicstyle=\small\ttfamily,
  keywordstyle=\color{magenta!75!black},
  stringstyle=\color{green!50!black},
  showspaces=false,
  showstringspaces=false,
  commentstyle=\color{gray}}

\title{
  \textbf{\large Projektaufgabe 3 } \\
  Phase 3 – Implementierung des XPath Accelerators (3.5 P) \\
  {\large Datenmanagement jenseits von Relationen}
}

\author{
  Gruppen Nummer (e.g. A1, B5, B3) \\
  \large Lastname1 Firstname1, StudentID1 \\
  \large Lastname2 Firstname2, StudentID2 
}

\begin{document}

\maketitle\thispagestyle{empty}

Dieses Reporting Template dient der Vorbereitung der Abgabe von Phase 3. 

\subsection*{Verkleinern der Fensteranfrage (0.5 P)}

Zeigen Sie, dass die Berechnung der Achsen für die Beispiele aus Phase 1 hier die gleichen Ergebnisse erzeugen:

\begin{table}[h]
	\centering
		\begin{center}
			\begin{tabular}{ l | c c }
				\toprule
				Achse & Ergebnisknoten ID's & Ergebnisgröße\\
				\midrule
				ancestor & \{0, 1, 16, 17\} & 4 \\ 
				descendants & $[1, 15] \cup [49, 61]$ & 28 \\  
				following SchmittKAMM23& \{49\} & 1 \\  
				preceeding SchmittKAMM23 & \{\} & 0 \\ 
				following SchalerHS23& \{\} & 0 \\  
				preceeding SchalerHS23& \{1\} & 1 \\ 
				\bottomrule
			\end{tabular}
			\end{center}
	\caption{Ergebnisse der XPath-Berechnung auf Edge Modell aus Phase 1}
		\label{tab:ErgebnisseDerXPathBerechnungEdgePhase3}
\end{table}

\begin{table}[h]
	\centering
		\begin{center}
			\begin{tabular}{ l | c c }
				\toprule
				Achse & Ergebnisknoten ID's & Ergebnisgröße\\
				\midrule
				ancestor & \{0, 1, 16, 17\} & 4 \\ 
				descendants & $[1, 15] \cup [49, 61]$ & 28 \\  
				following SchmittKAMM23& \{49\} & 1 \\  
				preceeding SchmittKAMM23 & \{\} & 0 \\ 
				following SchalerHS23& \{\} & 0 \\  
				preceeding SchalerHS23& \{1\} & 1 \\ 
				\bottomrule
			\end{tabular}
			\end{center}
	\caption{Ergebnisse der XPath-Berechnung des XPath-Accelerators aus Phase 3 mit Verkleinerung des Fensters}
		\label{tab:ErgebnisseDerXPathBerechnungReducedPhase3}
\end{table}



\subsection*{Schemaerstellung (1 Punkt).}
Geben Sie die Create Table statements aller von Ihnen angelegten Tabellen an:
\begin{lstlisting}[style=dmrsql]
CREATE TABLE accel (
    id BIGINT NOT NULL,
    pre BIGINT NOT NULL,
    post BIGINT NOT NULL,
    parent BIGINT
)

CREATE TABLE content (
    id BIGINT,
    text VARCHAR(512)
)

CREATE TABLE attribute (
    id BIGINT,
    text VARCHAR(512)
)
\end{lstlisting}

\subsection*{Zugriff mit nur einer Achse (1 Punkt).}
Zeigen Sie, dass die Berechnung der Achsen für die Beispiele aus Phase 1 hier die gleichen Ergebnisse erzeugen:

\begin{table}[h]
	\centering
		\begin{center}
			\begin{tabular}{ l | c c }
				\toprule
				Achse & Ergebnisknoten ID's & Ergebnisgröße\\
				\midrule
				descendants & $[1, 15] \cup [49, 61]$ & 28 \\  
				\bottomrule
			\end{tabular}
			\end{center}
	\caption{Ergebnisse der XPath-Berechnung des XPath-Accelerators aus Phase 3 mit nur einer Achse}
		\label{tab:ErgebnisseDerXPathBerechnungOneAxisPhase3}
\end{table}

Da sich die Knoten IDs unterscheiden, zeigen Sie für einige Knoten die jeweilige Identität.

\begin{table}[h]
	\centering
		\begin{center}
			\begin{tabular}{ l | l | l }
				\toprule
				Knoten ID & Identität & Typ\\
				\midrule
				17 & vldb\_2023 & publishingYear\\
				1 & SchmittKAMM23 & article\\
				49 & SchalerHS23 & article\\
				2 & Daniel Ulrich Schmitt & author\\
				\bottomrule
			\end{tabular}
		\end{center}
	\caption{Identität ausgewählter Knoten im Toy-Beispiel}
	\label{tab:KnotenIdentitaetenOneAxis}
\end{table}

\subsection*{Zeitmamagement}

Benötigte Zeit pro Person (nur Phase 2): \textbf{nicht angegeben}


\end{document}
